\scanpage{52}
% archive foliation: 51
% Sheet read after 180-degree rotation.
\hfill - 35 -\hfill n° 393

\deleted{
21) Soient A une algèbre de Banach unipollente, $\Delta$ une partie ouverte de $\mathbb{C}$, $\lambda\longrightarrow R(\lambda)$ une application de $\Delta$ dans A telle que:
\[
(1)\qquad R(\lambda)-R(\mu)=-(\lambda-\mu)R(\lambda)R(\mu)
\]
quels que soient $\lambda$, $\mu\in\Delta$.

a) Soit $\lambda_0\in\Delta$. Pour tout $\lambda\in\Delta$ tel que $|\lambda-\lambda_0|\,\|R(\lambda)\|<1$, on a:
\[
R(\lambda)=\sum_{n=0}^{\infty}(\lambda_0-\lambda)^n R(\lambda_0)^{n+1}.
\]

b) Pour tout $\lambda\in\Delta$, $\dfrac{d^n}{d\lambda^n}R(\lambda)=(-1)^n n!\,R(\lambda)^{n+1}$.

c) Si R est holomorphe à l'infini, il existe z, j, $x\in A$ tels que $z^2=0$, $j^2=j$, $zj=jz=0$, $x\in jAj$, et $R(\lambda)=z+jR(\lambda,x)$ pour $|\lambda|$ assez grand. (Ecrire $R(\lambda)=\sum_{n=0}^{\infty}c_n\lambda^{-n}$ pour $|\lambda|>\lambda_0$, et exprimer que R satisfait à (1). On peut poser $c_0=z$, $c_1=j$, $c_2=x$).

d) Pour que R soit la résolvante d'un élément de A, il faut et il suffit que $R(\lambda_0)$ soit inversible pour un $\lambda_0\in\Delta$.

e) Pour tout $x\in A$, la fonction $\lambda\longrightarrow\sum_{n=0}^{\infty}(\lambda_0-\lambda)^n x^{n+1}$ satisfait à (1) pour $|\lambda-\lambda_0|\,\|x\|<1$.

\P\quad 22) a) Soient B une algèbre de Banach unipollente, $w\in B$, $C>1$ et $\varepsilon>0$. Il existe $\eta>0$ tel que les conditions $u\in B$, $\|u\|<C$, $\|uw-w\|<\eta$, $0\leq\gamma\leq1/4C$ impliquent $\|yw-w\|<\varepsilon$, où $y=[(1-\gamma)e+\gamma u]^{-1}$.

b) Soit A une algèbre de Banach. On suppose qu'il existe une constante C telle que, pour tout $\varepsilon>0$ et tous $x_1,....,x_n\in A$, il existe $u\in A$ avec $\|u\|<C$, $\|ux_1-x_1\|<\varepsilon,....,\|ux_n-x_n\|<\varepsilon$. Pour tout $z\in A$, il existe $x,y\in A$ tels que $z=xy$, $y\in\overline{Az}$, $\|z-y\|<\delta$. (On peut supposer $C>1$. Soit $\gamma=1/4C$. Appliquant a) dans $\widetilde{A}$, déterminer par récurrence une suite $u_1,u_2,....$ dans A telle
}

